% GENERATED FILE. Edit canonical lessons, then run npm run generate:books.
% canonical-source-hash: fnv1a64:a5a7ae4ce535cee1
% canonical-lessons: SA-C410-chidram, SA-C410-ksarati, SA-C410-sithilah, SA-C410-vidirnah, SA-C410-gharsati

\chapter{Breakdowns}
\label{ch:sa-breakdowns}
\hlchaptermodality{410}

\begin{chapteropening}
\textbf{By the end of this chapter:} \emph{I can say what has broken or worn on a cart.}

Complete the last lesson of chapter 410: I can say what has broken or worn on a cart.
\end{chapteropening}

\section[\texorpdfstring{\sk{छिद्रम्} (chidram)}{छिद्रम् (chidram)}]{\sk{छिद्रम्} (chidram) — a hole}

\label{lesson:SA-C410-chidram}



\begin{quote}
Before the new one: say the Sanskrit for sets spinning, then the Sanskrit for a wheel track.
\end{quote}

\subsection*{You'll want to know: \sk{छिद्रम्}}
\textbf{\sk{छिद्रम्}} — \emph{chidram} — \textquotedblleft{}a hole\textquotedblright{}.

\begin{quote}
\sk{मम} \sk{स्यूते} \sk{छिद्रम्} \sk{अस्ति।}
\end{quote}

\begin{quote}
\emph{mama syūte chidram asti.}
\end{quote}

\begin{quote}
There is a hole in my bag.
\end{quote}

A neuter noun.

\begin{cousinweb}
\textbf{\sk{छिद्रम्}} (\emph{chidram}) comes from the root √chid, \textquotedblleft{}to cut, to split\textquotedblright{}, which also gives \emph{chinatti}, \textquotedblleft{}cuts\textquotedblright{}; MW gives \emph{chidra} as torn, pierced and, as a noun, a hole (MW).
\end{cousinweb}

\begin{grammarlens}[title={What you've built}]
One more word for getting around.
\end{grammarlens}

\subsection*{Your turn}
\begin{itemize}
\raggedright
  \item \emph{Say it:} \emph{chidram}
  \item \emph{Say it:} \emph{chidram}, once more
  \item \emph{Say it:} \emph{mama syūte chidram asti.} — There is a hole in my bag.
  \item \emph{Recall:} say the Sanskrit for sets spinning, then the Sanskrit for a wheel track, then say \emph{chidram} again
\end{itemize}

\begin{tcolorbox}[breakable,colback=teal!4,colframe=teal!35!black,title={Before you move on}]
What does \sk{छिद्रम्} mean? (\textquotedblleft{}A hole\textquotedblright{}.) Say it once more.
\end{tcolorbox}
\section[\texorpdfstring{\sk{क्षरति} (kṣarati)}{क्षरति (kṣarati)}]{\sk{क्षरति} (kṣarati) — leaks; trickles out}

\label{lesson:SA-C410-ksarati}



\begin{quote}
Before the new one: say the Sanskrit for a wheel track, then the Sanskrit for a hole.
\end{quote}

\subsection*{You'll want to know: \sk{क्षरति}}
\textbf{\sk{क्षरति}} — \emph{kṣarati} — \textquotedblleft{}leaks; trickles out\textquotedblright{}.

\begin{quote}
\sk{छिद्रात्} \sk{तैलं} \sk{क्षरति।}
\end{quote}

\begin{quote}
\emph{chidrāt tailaṃ kṣarati.}
\end{quote}

\begin{quote}
Oil leaks from the hole.
\end{quote}

Past (it …): \sk{तैलं} \sk{छिद्रात्} \sk{क्षरितवत्}, \emph{tailaṃ chidrāt kṣaritavat}.

\begin{cousinweb}
\textbf{\sk{क्षरति}} (\emph{kṣarati}) comes from the root √kṣar, \textquotedblleft{}to flow, to stream\textquotedblright{}; MW gives \emph{kṣarati} as to flow, trickle or drip, and KRM gives the participle \emph{kṣarita} (MW).
\end{cousinweb}

\begin{grammarlens}[title={What you've built}]
One more word for getting around.
\end{grammarlens}

\subsection*{Your turn}
\begin{itemize}
\raggedright
  \item \emph{Say it:} \emph{kṣarati}
  \item \emph{Say it:} \emph{kṣarati}, once more
  \item \emph{Say it:} \emph{chidrāt tailaṃ kṣarati.} — Oil leaks from the hole.
  \item \emph{Recall:} say the Sanskrit for a wheel track, then the Sanskrit for a hole, then say \emph{kṣarati} again
\end{itemize}

\begin{tcolorbox}[breakable,colback=teal!4,colframe=teal!35!black,title={Before you move on}]
What does \sk{क्षरति} mean? (\textquotedblleft{}Leaks; trickles out\textquotedblright{}.) Say it once more.
\end{tcolorbox}
\section[\texorpdfstring{\sk{शिथिलः} (śithilaḥ)}{शिथिलः (śithilaḥ)}]{\sk{शिथिलः} (śithilaḥ) — loose; slack}

\label{lesson:SA-C410-sithilah}



\begin{quote}
Before the new one: say the Sanskrit for a hole, then the Sanskrit for leaks.
\end{quote}

\subsection*{You'll want to know: \sk{शिथिलः}}
\textbf{\sk{शिथिलः}} — \emph{śithilaḥ} — \textquotedblleft{}loose; slack\textquotedblright{}.

\begin{quote}
\sk{अक्षस्य} \sk{कीलः} \sk{शिथिलः} \sk{अस्ति।}
\end{quote}

\begin{quote}
\emph{akṣasya kīlaḥ śithilaḥ asti.}
\end{quote}

\begin{quote}
The axle's pin is loose.
\end{quote}

\begin{cousinweb}
\textbf{\sk{शिथिलः}} (\emph{śithilaḥ}) comes from the root √śrath, \textquotedblleft{}to loosen\textquotedblright{}, by way of \emph{śithira}, \textquotedblleft{}loose\textquotedblright{}, of which MW gives it as a collateral form; MW gives \emph{śithila} as loose, slack, untied (MW).
\end{cousinweb}

\begin{grammarlens}[title={What you've built}]
One more word for getting around.
\end{grammarlens}

\subsection*{Your turn}
\begin{itemize}
\raggedright
  \item \emph{Say it:} \emph{śithilaḥ}
  \item \emph{Say it:} \emph{śithilaḥ}, once more
  \item \emph{Say it:} \emph{akṣasya kīlaḥ śithilaḥ asti.} — The axle's pin is loose.
  \item \emph{Recall:} say the Sanskrit for a hole, then the Sanskrit for leaks, then say \emph{śithilaḥ} again
\end{itemize}

\begin{tcolorbox}[breakable,colback=teal!4,colframe=teal!35!black,title={Before you move on}]
What does \sk{शिथिलः} mean? (\textquotedblleft{}Loose; slack\textquotedblright{}.) Say it once more.
\end{tcolorbox}
\section[\texorpdfstring{\sk{विदीर्णः} (vidīrṇaḥ)}{विदीर्णः (vidīrṇaḥ)}]{\sk{विदीर्णः} (vidīrṇaḥ) — split; burst open}

\label{lesson:SA-C410-vidirnah}



\begin{quote}
Before the new one: say the Sanskrit for leaks, then the Sanskrit for loose.
\end{quote}

\subsection*{You'll want to know: \sk{विदीर्णः}}
\textbf{\sk{विदीर्णः}} — \emph{vidīrṇaḥ} — \textquotedblleft{}split; burst open\textquotedblright{}.

\begin{quote}
\sk{शकटस्य} \sk{फलकं} \sk{विदीर्णम्} \sk{अस्ति।}
\end{quote}

\begin{quote}
\emph{śakaṭasya phalakaṃ vidīrṇam asti.}
\end{quote}

\begin{quote}
The cart's plank is split.
\end{quote}

\begin{cousinweb}
\textbf{\sk{विदीर्णः}} (\emph{vidīrṇaḥ}) is built from \emph{vi}, \textquotedblleft{}apart\textquotedblright{}, and the root √dṝ, \textquotedblleft{}to burst, to split\textquotedblright{}; MW gives \emph{vidīrṇa} as rent asunder, split, burst open (MW).
\end{cousinweb}

\begin{grammarlens}[title={What you've built}]
One more word for getting around.
\end{grammarlens}

\subsection*{Your turn}
\begin{itemize}
\raggedright
  \item \emph{Say it:} \emph{vidīrṇaḥ}
  \item \emph{Say it:} \emph{vidīrṇaḥ}, once more
  \item \emph{Say it:} \emph{śakaṭasya phalakaṃ vidīrṇam asti.} — The cart's plank is split.
  \item \emph{Recall:} say the Sanskrit for leaks, then the Sanskrit for loose, then say \emph{vidīrṇaḥ} again
\end{itemize}

\begin{tcolorbox}[breakable,colback=teal!4,colframe=teal!35!black,title={Before you move on}]
What does \sk{विदीर्णः} mean? (\textquotedblleft{}Split; burst open\textquotedblright{}.) Say it once more.
\end{tcolorbox}
\section[\texorpdfstring{\sk{घर्षति} (gharṣati)}{घर्षति (gharṣati)}]{\sk{घर्षति} (gharṣati) — rubs against; scrapes}

\label{lesson:SA-C410-gharsati}



\begin{quote}
Before the new one: say the Sanskrit for loose, then the Sanskrit for split open.
\end{quote}

\subsection*{You'll want to know: \sk{घर्षति}}
\textbf{\sk{घर्षति}} — \emph{gharṣati} — \textquotedblleft{}rubs against; scrapes\textquotedblright{}.

\begin{quote}
\sk{चक्रं} \sk{भित्तिं} \sk{घर्षति।}
\end{quote}

\begin{quote}
\emph{cakraṃ bhittiṃ gharṣati.}
\end{quote}

\begin{quote}
The wheel scrapes the wall.
\end{quote}

Past (it …): \sk{चक्रं} \sk{भित्तिं} \sk{घृष्टवत्}, \emph{cakraṃ bhittiṃ ghṛṣṭavat}.

\begin{cousinweb}
\textbf{\sk{घर्षति}} (\emph{gharṣati}) comes from the root √ghṛṣ, \textquotedblleft{}to rub, to grind\textquotedblright{}; MW gives the participle \emph{ghṛṣṭa}, \textquotedblleft{}rubbed, grazed\textquotedblright{}; it lives on as Hindi \emph{ghisnā}, \textquotedblleft{}to rub\textquotedblright{} (MW; CDIAL 4450).
\end{cousinweb}

\begin{grammarlens}[title={What you've built}]
One more word for getting around.
\end{grammarlens}

\subsection*{Your turn}
\begin{itemize}
\raggedright
  \item \emph{Say it:} \emph{gharṣati}
  \item \emph{Say it:} \emph{gharṣati}, once more
  \item \emph{Say it:} \emph{cakraṃ bhittiṃ gharṣati.} — The wheel scrapes the wall.
  \item \emph{Recall:} say the Sanskrit for loose, then the Sanskrit for split open, then say \emph{gharṣati} again
\end{itemize}

\begin{tcolorbox}[breakable,colback=teal!4,colframe=teal!35!black,title={Before you move on}]
What does \sk{घर्षति} mean? (\textquotedblleft{}Rubs against; scrapes\textquotedblright{}.) Say it once more.
\end{tcolorbox}
